\documentclass[10pt]{article}
\usepackage[margin=20mm]{geometry}
\usepackage{amsmath,amssymb,mathtools,fontspec,unicode-math,xurl,hyperref}
\IfFontExistsTF{DejaVu Serif}{\setmainfont{DejaVu Serif}}{\setmainfont{Latin Modern Roman}}
\setmathfont{latinmodern-math.otf}
\setlength{\parindent}{0pt}\setlength{\parskip}{6pt}\sloppy
\title{Formal initial review: M25-0078
The spectrum of excisive functors}
\author{Scientific review performed in this conversation}
\date{26 October 2026}
\begin{document}
\maketitle
\section*{Scope and initial outcome}
All 102 pages of the acquired publisher PDF were read, including the appendix and references. Locations below use one-based PDF page / printed page; printed page equals PDF page plus 362. Related author source arXiv:2402.04244v2 was parsed in full for structure and dependencies. It is not identified with the publisher’s TeX. All findings below concern the bound publisher PDF.\par
Initial outcome: Major. Two central proofs use an invalid intermediate inference: the passage from relative openness to ambient openness in Theorem 11.22, and the replacement of a proper-family Tate construction by ordinary cyclic Tate cohomology in Theorem 12.27. The first has a complete local repair. For the second, a direct elementary-abelian argument supplies the missing nonvanishing mechanism. Neither finding is a counterexample to the stated main classification. Four additional auxiliary assertions require restrictions or corrections.\par
No. | Confirmed assertion or inference | Judgment\par
1 | Theorem 11.22, 84/446: a finite relative open need not be ambient open. | Major\par
2 | Theorem 12.27, 96–97/458–459: the cyclic target is the wrong Tate construction. | Major\par
3 | Proposition 2.38, 18/380: arbitrary limit shapes do not give a smashing localization. | Minor\par
4 | Remark 2.28, 17/379: compact objects require stable, not only finite-colimit, generation. | Minor\par
5 | Corollary 8.13, 58/420: the first unital presentation omits \(x_1=1\). | Minor\par
6 | Definition 10.27 and Example 10.28, 76–77/438–439: non-\(p\)-local scope is inconsistent. | Minor\par
The severity reflects the role of the defective inference, not a claim that a repair is impossible. Finite computations supplement the proofs below and do not certify every imported theorem.\par
\section*{Problem 1: relative openness in the geometric blueshift proof}
Exact location. Theorem 11.22, paragraph following (11.25), 84/446; setup (11.24) on 83/445 and subsequent use on 84–85/446–447. Write \(X=\operatorname{Spc}(\mathrm{Exc}_d(\mathrm{Sp}^c,\mathrm{Sp})^c)\). The truncated spectrum is a finite subspace \(V=V_{\le h+n}\) of \(X\). The proof takes \[W=V\setminus(\overline{\{P_d([l],p,h)\}}\cap V)\] and asserts that its finiteness and relative openness make it open in \(X\), then chooses a compact \(z\) with \(\operatorname{supp}(z)=X\setminus W\). The second assertion requires ambient quasi-compact openness.\par
\subsection*{A legal instance where ambient openness fails}
Take \(d=2\), \(l=1\), \(p=2\), \(h=1\). The partition \(2\) has length one, so \(\delta_2(2,1)=1\) and the proof uses \(n=0\). Here \[V=\{P_2([1],0,1),P_2([2],0,1)\}.\] Proposition 11.7 excludes an inclusion between the two primes at equal finite height. Thus \(W=\{P_2([2],0,1)\}\).\par
Under the embedding \(\operatorname{Spc}(\partial_2):\operatorname{Spc}(\mathrm{Sp}^c)\to X\) from Proposition 7.7, the inverse image of this \(W\) is the singleton \(\{C_{0,1}\}\). That singleton is not open in the integral finite-spectrum Balmer spectrum. Indeed, a basic open \(U(a)\) containing \(C_{0,1}\) means that the finite spectrum \(a\) is rationally acyclic. Its integral homology is finite torsion in finitely many degrees, so only finitely many primes occur. For every other prime \(q\), \(a_{(q)}=0\). Consequently \(U(a)\) also contains the \(q\)-primary chromatic points. Every open neighborhood of \(C_{0,1}\) contains such a basic neighborhood, and cannot be the singleton. Continuity therefore proves that \(W\) is not open in \(X\). In particular \(X\setminus W\) cannot be the closed support claimed in the printed proof.\par
\subsection*{Complete repair of the support-selection step}
The actual argument only needs (11.25), not equality of the support with \(X\setminus W\). Since \(W\) is relatively open, choose an ambient open \(O\) with \(O\cap V=W\). For each \(w\in W\), choose a basic quasi-compact open \(B_w\) with \(w\in B_w\subset O\). The union \[U=\bigcup_{w\in W}B_w\] is a finite union of quasi-compact opens, and \(U\cap V=W\). Balmer’s Proposition 2.14 [Balmer] gives a compact \(z\) with \(U=X\setminus\operatorname{supp}(z)\). Hence \[\operatorname{supp}(z)\cap V=V\setminus W
 =\overline{\{P_d([l],p,h)\}}\cap V,\] exactly (11.25). If \(W\) is empty, take the tensor unit instead.\par
Downstream check and limits. The derivative tests in (11.27)–(11.29) only use this intersection identity, compactness, and dualizability. The triangle-inequality contradiction for \(\delta_p\) is unchanged. The Tate-vanishing step on the truncated category can use the right-adjoint identity in Remark 6.17. The support-splitting criterion used on 85/447 is [Retained latex: \textbackslash{}cite[Proposition 2.29]\{PSW\}], whose hypothesis is a compact object and a Thomason subset; those hypotheses hold after the finite localization. When \(h+n=1\), the relevant Tate derivatives vanish directly by Theorem 10.33, avoiding an unnecessary height-zero prime comparison. Thus the replacement repairs this specific gap. It is not a new independent proof of every earlier categorical construction.\par
Impact and counterevidence. This is a false inference in the core lower-bound proof, hence Major under the review standard. The support object actually needed does exist; the main theorem is not refuted. Working in an explicitly \(p\)-localized ambient category could also alter the openness issue, but the printed setup names the integral category and does not make that replacement.\par
\section*{Problem 2: ordinary and proper-family Tate constructions}
Exact location. Proof of Theorem 12.27, 96–97/458–459, the computation after (12.28) and the ring map on 97/459 attributed to Theorem 10.33. Compare the actual construction on 78/440, and Definitions 10.16–10.17.\par
For a nontrivial cyclic \(p\)-group \(C\), the restriction constructed in Theorem 10.33 has target \[(H\mathbb F_p)^{t\mathcal P(C)}
 =\Phi^C\bigl(\text{Borel-completed }H\mathbb F_p\bigr),\] where \(\mathcal P(C)\) is the family of all proper subgroups. The proof of Theorem 12.27 instead writes \((H\mathbb F_p)^{tC}\) and asserts nonvanishing using ordinary Tate cohomology. The two constructions coincide for \(C=C_p\), but not for higher prime-power order. The distinction is explicit in the paper’s own family notation and in [Retained latex: \textbackslash{}cite[Definition 3.3]\{BHN\}].\par
\subsection*{The actual target is zero for a permitted cyclic group}
Take \(k=d=4\), \(l=1\), \(p=2\), \(\lambda=(4)\). The long-cycle subgroup is \(C=C_4\subset\Sigma_4\). Its proper subgroups are \(e\) and \(C_2\), both nontransitive, so the restricted family is precisely \(\mathcal P(C_4)\).\par
The reduced real regular representation is a sign line plus a real rotation plane: \(\bar\rho=\epsilon\oplus\eta\). Let \(v\in H^1(BC_4;\mathbb F_2)\) be the parity homomorphism. It satisfies \(v^2=0\). This can be checked without a cohomology-ring lookup: for representatives \(a=0,1,2,3\), put \(b(a)=\lfloor a/2\rfloor\bmod2\). With trivial coefficients, \[\delta b(a,c)=b(c)-b(a+c\bmod4)+b(a)
             =(a\bmod2)(c\bmod2).\] All sixteen pairs were also checked exactly in the independent script. Thus the mod-two Euler class of \(\bar\rho\) satisfies \[e(\bar\rho)=v\,w_2(\eta),\qquad e(\bar\rho)^2=0.\] The representation-sphere model of Example 10.10 and Lemma 10.24, together with the mod-two Thom isomorphism, identifies the homotopy groups of the proper-family construction with the localization of \(H^*(BC_4;\mathbb F_2)\) at this Euler class. Equivalently, its direct-system maps multiply by \(e(\bar\rho)\). Every element becomes zero after at most two such maps. Therefore \[(H\mathbb F_2)^{t\mathcal P(C_4)}=0.\] In contrast, ordinary \((H\mathbb F_2)^{tC_4}\) is nonzero: the complete periodic cyclic resolution has alternating maps \(g-1\) and \(1+g+g^2+g^3\), both zero on the trivial \(\mathbb F_2\) module. The asserted nonzero target cannot be the target constructed earlier.\par
What this does and does not disprove. It disproves the cited nonvanishing argument for a legal partition. It does not show that the \(\Sigma_4\) family factor vanishes, or that the prime inclusion in Theorem 12.27 is false. Indeed, restriction to the regular subgroup \(C_2^2\) gives the Euler class \(xy(x+y)\) in \(\mathbb F_2[x,y]\), which is nonnilpotent. That provides the required nonzero target in this instance.\par
\subsection*{A repair using elementary abelian groups}
For any \(p\)-power partition \(\lambda=(p^{e_1},\ldots,p^{e_l})\) with \(l<k\), use \[E=\prod_{i=1}^l E_i,\qquad E_i=(C_p)^{e_i},\] with the regular permutation action of \(E_i\) on its block. The restriction of the representation in Lemma 10.24 is \(V=\bigoplus_i\bar\rho_{E_i}\), inflated to the product. Its Euler class is a product of nonzero linear forms in the polynomial part of \(H^*(BE;\mathbb F_p)\), up to nonzero scalar factors. For \(p=2\) these are the nontrivial real characters in degree one. For odd \(p\), the nontrivial characters pair into oriented real planes; their Euler classes are nonzero degree-two linear forms. Blocks of size one contribute the unit.\par
The polynomial part is an integral domain. For odd \(p\), the full cohomology is free over that polynomial part with an exterior basis. Consequently \(e(V)\) is not nilpotent and \[H^*(BE;\mathbb F_p)[e(V)^{-1}]\ne0.\] The representation-sphere model and restriction give a unital ring map from \((H\mathbb F_p)^{t\mathcal F_{\rm nt}(\lambda)}\) to this nonzero construction. It follows that the factor is nonzero. Proposition 10.26 then gives nonzero \(\partial_l t_k i_k(H\mathbb F_p)\). This is a general argument; the finite polynomial calculations are only supporting checks.\par
\subsection*{Coefficient exchange should also be made explicit}
The printed calculation on 96/458 invokes Proposition 6.14 to commute free \(H\mathbb Z\) extension with a Tate construction. The Tate assertion there has a closed-monoidal hypothesis; it is not justified merely by a strong monoidal left adjoint. This review does not count a separate counterexample to that particular exchange.\par
One can avoid the unestablished exchange. First pass to the \(k\)-excisive truncation, which preserves compactness and the derivatives of degrees at most \(k\). Let \(U\) be the forgetful functor from \(H\mathbb Z\)-module-valued functors to spectrum-valued functors. It preserves limits and colimits, hence commutes with the derivative construction. The right-adjoint formula of Remark 6.17, applied to \(t=i_k(S/p)\), identifies \[U\bigl(t_k^\mathbb Zi_k^\mathbb Z(H\mathbb F_p)\bigr)
       \simeq t_k i_k(H\mathbb F_p).\] Here free extension sends \(i_k(S/p)\) to \(i_k^\mathbb Z(H\mathbb F_p)\), and \(t\otimes U1\) is \(i_k(H\mathbb F_p)\). Thus \(M=\partial_l^\mathbb Zt_k^\mathbb Zi_k^\mathbb Z(H\mathbb F_p)\) is a nonzero \(H\mathbb F_p\)-module by the elementary-abelian argument.\par
For the compact object \(x\) in (12.28), dualizability and \(\partial_k^\mathbb Z(i_k^\mathbb Z(H\mathbb F_p)\mathbin{\circledast}x)=0\) give \[0=M\otimes_{H\mathbb Z}\partial_l^\mathbb Zx
   \simeq
 M\otimes_{H\mathbb F_p}
       \bigl(H\mathbb F_p\otimes_{H\mathbb Z}\partial_l^\mathbb Zx\bigr).\] Tensor product with a nonzero complex over the field \(\mathbb F_p\) detects zero: choose nonzero homology classes in both factors and use the field Künneth formula. Hence \(H\mathbb F_p\otimes_{H\mathbb Z}\partial_l^\mathbb Zx=0\). This proves the required inclusion for a single \(p\)-power partition; the partition-chain argument handles the remaining congruent pairs.\par
Impact and repair limits. The coefficient-spectrum classification has a Major proof gap as printed. The above replacement gives its missing algebraic nonvanishing and uses an explicit right-adjoint comparison. The full coherence of the surrounding Day-convolution and module category constructions has not been reconstructed from first principles here; no scientific acceptance is inferred from this repair.\par
\section*{Problem 3: arbitrary limit shapes in Proposition 2.38}
Location. Proposition 2.38 and its proof, 18/380, with the unrestricted small family of diagrams in Remark 2.37. The argument uses left exactness of filtered colimits to assert closure of the local objects under filtered colimits. Left exactness only controls finite limit shapes.\par
\subsection*{Counterexample satisfying the stated tensor stability}
Fix a prime \(p\). For every finite spectrum \(X\), include in \(J\) the inverse tower \[\cdots\xrightarrow{p}X\xrightarrow{p}X\] with cone vertex \(0\) and zero cone maps. Use an essentially small skeleton and its tensor closure. This is a small set of diagrams, closed under postcomposition with tensoring by any finite spectrum.\par
For \(j\ge1\) let \(F_j(X)=H(\mathbb Z/p^j)\otimes X\). These are reduced, linear functors. On \(H(\mathbb Z/p^j)\) the \(j\)-fold composite in the inverse \(p\)-tower is zero. Its inverse limit and derived inverse limit vanish; explicitly the tower is pro-zero, and the Milnor sequence gives a contractible homotopy limit. Tensoring with a finite \(X\) preserves that limit because \(X\) is dualizable. Therefore every \(F_j\) sends every diagram in \(J\) to a limit diagram.\par
The maps \(\mathbb Z/p^j\to\mathbb Z/p^{j+1}\), \(a\mapsto pa\), induce a filtered system with colimit \[F(X)=H(P)\otimes X,\qquad P=\mathbb Z[1/p]/\mathbb Z.\] At \(X=S\), the inverse \(p\)-tower on \(P\) has a nonzero compatible element \[x_n=p^{-(n+1)}+\mathbb Z,\qquad p x_{n+1}=x_n.\] Because \(\pi_1 HP=0\), the Milnor sequence identifies \(\pi_0\operatorname{holim}(\cdots HP\xrightarrow p HP)\) with \(\lim_p P\), which contains this nonzero element. Thus \(F\) is not \(J\)-local.\par
Consequence. The inclusion of local objects does not preserve this filtered colimit, so the localization endofunctor is not smashing in Definition 2.34. This does not challenge the existence of the localization or tensor compatibility alone.\par
Repair and counterevidence. Restrict the allowed shapes to finite limit diagrams, in particular finite poset cubes. Colimits in spectra commute with finite limits; the local subcategory is then colimit closed. The actual Goodwillie-excision application in Theorem 2.39 uses finite cubes, so that application survives this restriction. The finding concerns auxiliary generality and is Minor. An assumption ensuring colimit closure by another argument would also suffice; tensor stability of \(J\) alone does not.\par
\section*{Problem 4: finite colimits do not generate all compact objects}
Location. Remark 2.28, 17/379. It characterizes compact functors as retracts of finite colimits of corepresentable functors \(h_X=\Sigma^\infty\operatorname{Map}(X,-)\).\par
Each \(h_X\) is objectwise connective. Spectra that are connective stay connective under finite colimits and retracts, and these operations in the functor category are objectwise. But \(\Sigma^{-1}h_S\) is compact: corepresentables are compact and suspension is an autoequivalence preserving compactness. At \(S\) its \(\pi_{-1}\) is \(\pi_0 h_S(S)\ne0\), since the stable degree components of \(\operatorname{Map}(S,S)\) give nontrivial reduced zeroth homology. It cannot lie in the stated finite-colimit/retract closure.\par
Repair and limits. Use the thick stable closure of the corepresentables, or finite colimits and retracts of all their integer suspensions. Compact generation, dualizability, and closure under the objectwise smash product then extend through this stable closure. The cited Ching Lemma 4.14 was checked and uses the same abbreviated finite-colimit wording; this report does not attribute a missing qualification to that reference or claim the wording originated here. The explicit closure characterization needs the correction, but no main classification is refuted. Judgment: Minor.\par
\section*{Problem 5: the missing unit relation in Corollary 8.13}
Location. First presentation and proof of Corollary 8.13, 58/420; relations (8.11) on 57/419.\par
Let \(B=\mathbb Z[x_1,\ldots,x_d]/I\), with \(I\) generated only by the displayed multiplication relations \[x_i x_j=\sum_l\mu(i,j,l)x_l .\] Evaluation at all \(x_i=0\) is a unital map \(B\to\mathbb Z\), so these relations cannot imply \(x_1=1\). Already for \(d=1\), \[B=\mathbb Z[x_1]/(x_1^2-x_1)\cong\mathbb Z\times\mathbb Z,
 \qquad A(1)=\mathbb Z.\] This is a counterexample to the first presentation, not merely a failure of its proof.\par
More generally the product rules give \(x_1x_i=x_i\), and all positive degree monomials reduce to integer combinations of the \(x_i\). There is an explicit isomorphism \[B\longrightarrow \mathbb Z\times A(d),\qquad
 1\longmapsto(1,1),\quad x_i\longmapsto(0,x_i).\] Surjectivity follows using \(1-x_1\) and the \(x_i\). For injectivity reduce any polynomial to \(a_0+\sum_i a_i x_i\) and use the given free abelian basis of \(A(d)\): the first coordinate determines \(a_0\), and then the second determines all \(a_i\). Thus the omitted relation adds a whole extra factor.\par
Repair and counterevidence. Add \(x_1-1\) to \(I\), or use the paper’s second presentation over \(\mathbb Z[x_2,\ldots,x_d]\) with \(x_1\) interpreted as \(1\). The actual construction of \(A(d)\) by its mark embedding and the identification \(R(d)\cong A(d)\) through the derivative marks remain valid. The separate malformed product display in Theorem 9.25 is an evident index/typesetting error whose intended rule is fixed by its proof; it is not counted as another finding. Judgment: Minor.\par
\section*{Problem 6: nonlocal spectra in the height convention}
Locations. Definition 10.27, Example 10.28 and Remark 10.29, 76–77/438–439; \(C_{p,0}=\mathrm{Sp}^c_{(p)}\) in Example 6.5, 41/403.\par
The definition takes the infimum over \(h\in\mathbb N\cup\{-1,\infty\}\) such that \(M\otimes C_{p,h+1}=0\). Example 10.28 then asserts that every spectrum which is not \(p\)-local has infinite \(p\)-height. Choose a distinct prime \(q\) and \(M=H\mathbb F_q\). For every finite \(p\)-local spectrum \(X\), multiplication by \(q\) is invertible on \(X\) and zero on \(H\mathbb F_q\). Therefore \(H\mathbb F_q\otimes X=0\). The \(h=-1\) condition holds, and the literal definition gives \[\operatorname{height}_p(H\mathbb F_q)=-1.\] Yet this nonzero spectrum is not \(p\)-local. It also cannot satisfy \(M\simeq L^f_{p,h}M\) for any listed \(h\): at \(h=-1\) the target is zero, and for other \(h\) it is \(p\)-local. The rewritten infimum in Remark 10.29 therefore gives infinity.\par
Repair and limits. Restrict the definition and equivalence to \(p\)-local \(M\), replacing the sphere example by \(S_{(p)}\), or explicitly define the height of arbitrary \(M\) through \(M_{(p)}\) and rewrite the subsequent criterion accordingly. The spectra \(L^f_{p,h}S\) and their Tate derivatives used in Theorem 10.33 are \(p\)-local, so this counterexample does not refute that theorem. It concerns the advertised scope of the definition and is Minor.\par
\section*{Positive checks, proof dependencies, and coverage}
\subsection*{Definitions, monoidal structures, and generators}
The functors are reduced; Day convolution is distinguished from objectwise smash throughout this review. The source is finite spectra, with finitary extension where used. The \(d\)-excisive localization and the finite derivative family were tracked through the recollement, compact generation, and prime classification. The Boolean cross-effect idempotents were checked algebraically: with \(\psi_U\psi_V=\psi_{U\cap V}\), \[c_U=\sum_{V\subset U}(-1)^{|U|-|V|}\psi_V\] are mutually orthogonal and sum to the identity. Exact finite checks cover up to five inputs. The empty-subset convention in one product rewrite does not alter the first or third defining formula or the subsequent arguments; it is not counted as a substantive finding.\par
For the generators \(P_dh_S(i)\), the derivative formula uses \(|\operatorname{surj}(k,i)|\) copies of the sphere, with the appropriate vanishing for \(k<i\). The tensor product formula counts subsets of \([i]\times[j]\) projecting onto both factors. Independent enumeration for \(i,j\le3\) agrees with the inclusion-exclusion formula. The surjection-product identity was checked for \(k\le12\). For the integral mark matrix, \[M_{ki}=i!\,\left\{\begin{matrix}k\\i\end{matrix}\right\},
 \qquad M=L\,\operatorname{diag}(1!,\ldots,d!),\] the Stirling matrix \(L\) is lower triangular and unimodular. This justifies the cokernel factors \(\mathbb Z/i!\), rather than inferring them from a triangular diagonal alone. Smith forms were checked through \(d=8\). Mark-prime congruences were checked for \(d=20\) and primes at most 19. The trace generators in Section 9 and the \(d!\) transfer argument were checked against this mark description.\par
\subsection*{Tate derivatives, height, and partition chains}
Proposition 10.26 was followed through its partition indexing, Young subgroup, and representation-sphere family. The vanishing direction uses a Sylow subgroup contained in the nontransitive family and a prime-to-\(p\) index. For Theorem 10.33, ordinary Tate vanishing supplies the upper bound. For the lower bound the diagonal cyclic subgroup has order the largest block size. Every proper subgroup fails transitivity on a block of maximal size; this is the needed reason the restricted family consists of all proper subgroups. It need not fail transitivity on every smaller block. The abelian-group height theorem cited there concerns the finite-local sphere and gives a drop by the cyclic group’s \(p\)-rank one [Retained latex: \textbackslash{}cite[Theorem 1.5]\{BHN\}]. It does not justify substituting ordinary \(H\mathbb F_p\) Tate cohomology in Problem 2.\par
The base-\(p\) digit sum, splitting/carrying operations and the two-step partition construction were checked analytically. Dynamic programming independently enumerated partition lengths and shortest chains for \(1\le l\le k\le128\) at \(p=2,3,5,7,11\): 41,280 pairs, including 720 two-step cases. They agree with \[\delta_p(k,l)=
 \begin{cases}
 0 & k=l,\\
 1 & k>l,\ k\equiv l\pmod{p-1},\ l\ge s_p(k),\\
 2 & k>l,\ k\equiv l\pmod{p-1},\ l<s_p(k),\\
 \infty & \text{otherwise}.
 \end{cases}\] This finite enumeration is not a substitute for the general proof. The height shifts, the distinction between equal finite heights and infinite height, and the restriction \(p=q\) away from height one were checked in the prime-inclusion statement.\par
\subsection*{Applications and boundaries of the audit}
Thomason subsets, admissible type functions, and the order of the correspondence were traced through Section 12. The type uses nonvanishing Morava \(K\)-homology; the original page image confirms the nonzero symbol that extraction can lose. The Smith/Floyd discussion uses the appropriate compactness, idempotent actions, and field dimensions; the separate caveat on multiplicative Morava theories at the prime two was read. The coefficient classification, noetherian description, and stratification/costratification/telescope applications were followed through their stated dependencies. Appendix A and all references were read.\par
Not fully reconstructed are all infinity-categorical coherences of Day convolution, Kan extension and module comparisons; full external nilpotence/chromatic and stratification theorems; and every detail of the external Smith-to-Floyd idempotent-filtration argument. Those limitations are recorded as remaining scientific work, not converted into confirmed defects. Isolated index or extraction errors with an unambiguous correct argument were excluded.\par
\section*{Supporting checks}
No. | Judgment | Proposed repair | Remaining limit\par
1 | Major | Extend the relative open by a finite union of ambient basic opens. | Does not reprove preceding categorical results.\par
2 | Major | Use regular elementary abelian subgroups and the right-adjoint Tate identity. | Full categorical framework remains an imported dependency.\par
3 | Minor | Restrict to finite limit shapes or assume colimit closure. | Does not repair arbitrary infinite-shape claims.\par
4 | Minor | Include stable closure and all integer suspensions. | No full reconstruction of cited compact-generation theorem.\par
5 | Minor | Impose \(x_1=1\). | Other ring identification unchanged by this repair.\par
6 | Minor | Restrict to \(p\)-local spectra or consistently localize \(M\). | Global-height wording must be changed.\par
No author adoption was sought. Every entry is an initial-review judgment, and `human\_reviews=[]`, `scientific\_decision=pending` remain unchanged.\par
The primary PDF SHA-256 is\par
eeebcaefa344432bf83172e6c22989b1
e702f675ddaac49bde3673d68ecac7b0\par
Its size is 2,365,377 bytes. The 13 manifest entries were fully parsed and checked at intake. The author source includes `goodwillie.tex` (302,072 bytes, 4,047 lines), `preamble.tex` (20,237 bytes, 637 lines), `quiver.sty` (1,606 bytes, 40 lines), the bibliography and diagram PDFs. Static parsing is not counted as scientific reading. No author program or author TeX was executed. Absence of computational code is not a defect.\par
This exact report source is compiled using the existing authorized Tectonic 0.17.0 in untrusted mode. The cache snapshot is not mislabeled as an exact trace of used files. The final report is released only after every report-page image has actually been viewed; the page-check record distinguishes all 102 text pages read from the 54 source-page images inspected.\par
No independent backend equivalence is asserted. All writes are confined to this task’s staging directory. Artifact completeness does not constitute scientific acceptance.\par
9 G. Arone, T. Barthel, D. Heard, and B. Sanders, The spectrum of excisive functors, Invent. Math. 241 (2025), 363–464. Bound publisher PDF; page-specific evidence given above. P. Balmer, The spectrum of prime ideals in tensor triangulated categories, author manuscript dated December 11, 2004, Proposition 2.14, PDF p. 7. https://www.math.ucla.edu/\textasciitilde{}balmer/Pubfile/Spectrum.pdf. I. Patchkoria, B. Sanders, and C. Wimmer, The spectrum of derived Mackey functors, author manuscript dated April 27, 2021, Proposition 2.29, PDF pp. 12–13. https://people.ucsc.edu/\textasciitilde{}beren/pdfs/derived-mackey.pdf. T. Barthel, M. Hausmann, N. Naumann, T. Nikolaus, J. Noel, and N. Stapleton, The Balmer spectrum of the equivariant homotopy category of a finite abelian group, arXiv:1709.04828v3, December 7, 2018. Theorem 1.5 and Definition 3.3. https://arxiv.org/pdf/1709.04828. M. Ching, Infinity-operads and Day convolution in Goodwillie calculus, J. London Math. Soc. (2) 104 (2021), 1204–1249. Publisher PDF read online, Lemma 4.14, PDF pp. 19–20. https://www.sas.rochester.edu/mth/sites/doug-ravenel/otherpapers/Ching-Day-conv.pdf.\par
\end{document}
